
\documentclass[twocolumn,10pt]{article}
%\usepackage{times}
%\documentclass[12pt,thmsa]{article}
\usepackage{amsfonts}
%%%%%
\newtheorem{theorem}{Theorem}
\newtheorem{corollary}[theorem]{Corollary}
\newtheorem{conjecture}[theorem]{Conjecture}
\newtheorem{lemma}[theorem]{Lemma}
\newtheorem{proposition}[theorem]{Proposition}
\newtheorem{definition}[theorem]{Definition}
\newtheorem{example}[theorem]{Example}
\newtheorem{axiom}{Axiom}
\newtheorem{remark}{Remark}
\newtheorem{exercise}{Exercise}[section]
\def\qqq{\bf{Q}}
\def\ccc{\bf{C}}
\def\zzz{\bf{Z}}
\def\rrr{\bf{R}}
\def\aaa{\bf{A}}
\def\nnn{\bf{N}}

\begin{document}
\author{D. Joyner}

\title{On the unitary dual of 
$\overline{SL(2,F)}$, $F$ p-adic}

\date{7-21-99}
\maketitle

%\begin{abstract}
\begin{quotation}
{\bf Abstract}: This paper describes the genuine irreducible unitary
representations of the $n$-fold metaplectic cover
$\overline{SL(2,F)}$, $F$ p-adic, 
with the exception of the reducible principal series. 
The image of the metaplectic Harish-Chandra
transform is determined and then a transfer of 
smooth (resp., spherical)  genuine functions
on $\overline{SL(2,F)}$ with their
analogs on $SL(2,F)$ is given such that their
associated orbital integrals are equal.
Finally, it is noted that the corresponding
principal series representations on $SL(2,F)$ and on
$\overline{SL(2,F)}$ satisfy a simple character relation.
%\end{abstract}
\end{quotation}

\tableofcontents
\vskip .4in

\section{Introduction}
Some of the results found in this paper:
\begin{itemize}
\item
Using the theory of germs of orbital integrals on a metaplectic group
\cite{V}, one can classify the image of the metaplectic Harish-Chandra
transform (see Proposition \ref{prop:main1}
and Corollary \ref{cor:main1}).
Moreover, it is not
hard to describe a transfer (or ``matching'') of 
smooth (resp., spherical)  genuine functions
on $\overline{SL(2,F)}$ with their
analogs on $SL(2,F)$ such that the
associated orbital integrals are equal.

\item
Based on the work of Moen \cite{Mo1}, \cite{Mo2},
Ariturk \cite{Ar}, Kazhdan-Patterson \cite{KP},
and Flicker-Kazhdan \cite{FK},
one can describe the genuine irreducible unitary
representations of the $n$-fold metaplectic cover
$\overline{SL(2,F)}$,
with the exception of the reducible principal series, 
where $F$ is a p-adic field (see for example
Propositions \ref{prop:moen} and \ref{prop:fk}).
In particular, some of the main results of Moen and 
Ariturk can be extended to the case where 
$n\geq 1$ is arbitrary. We describe
the special representation of $\overline{SL(2,F)}$
in \S \ref{subsec:special}. The supercuspidal representations
have already been described - see, 
for example, \cite{J3} (also \cite{M1}, \cite{M2}, \cite{J2}).

\item
Finally, assuming $gcd(p,n)=1$, we note that the 
corresponding principal
series representations on $SL(2,F)$ and its $n$-fold cover
(see Theorems \ref{thrm:main1} and \ref{thrm:main2})
satisfy simple character relations of the usual sort.
Under this correspondence, unitary (resp., spherical) representations match
with unitary (resp., spherical) representations. 
\end{itemize}

Y. Flicker \cite{F} described a comparison of smooth functions
on $G=GL(2,F)$ and an $n$-fold metaplectic
cover $\overline{G}$ by matching orbital integrals.
Flicker-Kazhdan \cite{FK} (resp., the author \cite{J1})
determined a comparison of smooth functions
on $G=GL(n,F)$ and on $\overline{G}$ 
(resp., on Cartans $T$ of $\overline{SL(n,F)}$ 
and on $\overline{SL(n,F)}$ itself) using the matching of
orbital integrals. In each of these papers, the authors then
obtained an analogous correspondence of
certain representations using character theory.
The idea here is similar but to instead compare orbital
integrals, and principal series representations,
on $G=SL(2,F)$ with those on $\overline{G}$.

These results, in the case $n=2$, form a 
necessary step in what will hopefully be a
proof of Labesse's multiplicity one conjecture for
$SL(2,{\Bbb{A}})$ (in preparation). 

{\bf Acknowledgements}:
 I am grateful to Courtney Moen for his helpful criticisms
on this paper. Also, I thank an anonymous referee for
several useful suggestions and corrections.

\section{Background notation}

Let $F$ be a p-adic field of characteristic 0, 
with uniformizer $\pi$, ring of integers ${\cal O}={\cal O}_F$, and let 
$q$ denote the cardinality of the residue field.
Let ${\Bbb{C}}$ denote the field of complex numbers and
${\Bbb{Z}}$ the ring of integers.
Let $\mu_n(F)$ denote the group of $n^{th}$ roots of unity
in $F$ and assume $|\mu_n(F)|=n$. We 
will identify $\mu_n(F)$ with a subgroup 
of ${\Bbb{C}}^\times$ (via some fixed isomorphism 
$\theta:\mu_n(F)\rightarrow \mu_n({\Bbb{C}})$).
Write $\mu_n=\mu_n(F)$.
Let $G=SL(2,F)$. Let $\overline{G}=\overline{SL(2,F)}$ 
denote the $n$-fold metaplectic cover defined by the cocycle
\begin{equation}
\beta_0(g_1,g_2)=({x(g_1g_2)\over x(g_1)},{x(g_1g_2)\over x(g_2)}),
\label{eqn:1}
\end{equation}
for $g_1,g_2\in G$, where
\[
x(
\begin{array}{cc}
a & b\\
c & d
\end{array}
)=
\left\{
\begin{array}{cc}
c,& {\rm if}\ c\not= 0,\\
d,& {\rm otherwise},
\end{array}
\right.
\]
and where $(\ ,\ )=(\ ,\ )_n$ is the $n^{th}$ power Hilbert symbol. 
By \S 1 of Kubota \cite{K}, $\beta_0$ is a 2-cocycle 
(called a ``factor set'' in \cite{K}) which is trivial
near the identity in $G\times G$, so $\overline{G}$ is a 
topological covering group of $G$. We denote
elements of $\overline{G}$ by $\overline{x}=(x,\zeta)$,
$x\in G,\ \zeta\in \mu_n$, and multiplication by
\[
\overline{x}_1\overline{x}_2=(x_1,\zeta_1)(x_2,\zeta_2)
=(x_1x_2,\beta_0(x_1,x_2)\zeta_1\zeta_2).
\]
Let $\rho:\overline{G}\rightarrow G$ denote
the natural surjection. If $H$ is any subset
of $G$ then we define $\overline{H}=\rho^{-1}(H)$.
Let $K_0=SL(2,{\cal O})$. If $p$ is
relatively prime to $2n$ then it is known
(see \cite{K}) that $\overline{K_0}$ splits,
i.e., that there is an isomorphism
$\overline{K_0}\cong K_0\times \mu_n$, which we
fix once and for all. In this case, we sometimes
identify $K_0$ with a subgroup of $\overline{K_0}$
via this isomorphism.

We choose a Haar measure on $G$, to be normalized
below, and a Haar measure on $\mu_n$ normalized so that 
$vol(\mu_n)=1$. Let the Haar measure on $\overline{G}$ be
the product measure.

For any subset $H\subset G$, let $H_{reg}$
denote the collection of elements of
$H$ which have distinct eigenvalues (in
some algebraic extension of $F$).

Let $A_0$ denote the diagonal subgroup of
$G$, let $W=N_G(A_0)/A_0$ denote the Weyl
group, represented by the matrices
\[
I=
\left(
\begin{array}{cc}
1 & 0 \\
0 & 1
\end{array}
\right),\ \ \ \ 
w_0=
\left(
\begin{array}{cc}
0 & 1 \\
-1 & 0
\end{array}
\right).
\]
Here $N_G$ denotes the normalizer in $G$ as usual.

Let $N_0$ denote the subgroup of 
$G$ consisting of upper triangular unipotent matrices 
and let $B=A_0N_0=N_0 A_0$ denote the upper
triangular (standard) Borel subgroup.
It is easy to verify that $\overline{N_0}$ splits.
We often identify $N_0$ with the
subgroup $(N_0,1)$ of $\overline{N_0}$.

\begin{definition}
Recall that a complex representation of $\overline{G}$
on a vector space $V$,
$\overline{\pi}:\overline{G}\rightarrow {\rm Aut}(V)$,
is called {\bf genuine} if $\overline{\pi}(g,\zeta)
=\theta(\zeta)\overline{\pi}(g,1)$, for
$g\in G$ and $\zeta\in \mu_n$.

Recall that a complex-valued function on $\overline{G}$
$f:\overline{G}\rightarrow {\Bbb{C}}$,
is called {\bf genuine} if $f(g,\zeta)
=\theta(\zeta)^{-1}f(g,1)$, for
$g\in G$ and $\zeta\in \mu_n$.
\end{definition}

We shall omit the use of $\theta$ and simply
identify $\mu_n$ with $\mu_n({\Bbb{C}})$.

Here we call an irreducible 
admissible represention of $G$ (or $\overline{G}$ 
{\bf supercuspidal}
if all its matrix coefficients are compactly supported.

\section{Orbital integrals}


Let
$C_c^\infty(\overline{G})_{gen}$
denote the algebra of compactly supported
locally constant genuine 
(i.e., ``smooth'')
functions on $\overline{G}$ and let
${\cal H}(\overline{G}//K_0)_{gen}$
denote the algebra of compactly supported
locally constant genuine bi-$K_0$-invariant
(i.e., ``spherical'')
functions on $\overline{G}$.
Let $M$ denote a Levi subgroup of a parabolic
subgroup $P=MN$, let
\[
\Delta(\overline{m})=\Delta_M(\overline{m})
=|\det(Ad_{\frak{n}}(m)-1_{\frak{n}})|,
\]
and let
\[
\delta(\overline{m})=\delta_M(\overline{m})
=|\det\ Ad_{\frak{n}}(m)|,
\]
where $\overline{m}\in \overline{M},\ 
\overline{m}=(m,\zeta)$, and where
$\frak{n}$ denotes the Lie algebra of $N$.
If 
$m=\left( \begin{array}{cc}
a & 0 \\
0 & a^{-1} \end{array}
\right)$ and $M$ is the diagonal subgroup of $G$ then
$\Delta(\overline{m})=|1-a^{-2}|$ and
$\delta(\overline{m})=|a^2|$.

Let $A=A_0$ from this point on, unless stated otherwise.
Let $S^{\overline{A}}$ denote the orbital integral
(restricted to a certain subset of 
$\overline{A}$),
\[
\begin{array}{c}
S^{\overline{A}}f(\overline{m})=\Delta(m)\delta(m)^{-1/2}
\int_{\overline{G}/\overline{A}}
f(g\overline{m}g^{-1})\ dg \\
=\delta(m)^{1/2}\int_{N}f(n\overline{m})\ dn,
\end{array}
\]
provided $Cent(\overline{m},\overline{G})=\overline{A}$.
We will show that $S$ defines algebra homomorphisms
\[
C_c^\infty(\overline{G})_{gen}
\rightarrow 
C_c^\infty(\overline{A}^{n_0})^W_{gen}
\]
and 
\[
{\cal H}(\overline{G}//K_0)_{gen}
\rightarrow 
{\cal H}(\overline{A}^{n_0}//K_0\cap A^{n_0})^W_{gen}.
\]
{\it Here, and throughout the rest of this paper, let}
\[
n_0=n/gcd(n,2).
\]
Moreover, the subalgebra of $W$-invariants,
\[
C_c^\infty(\overline{A})^W_{gen}
=\{f\in C_c^\infty(\overline{A}^{n_0})_{gen} \ |\ 
f((w,1)^{-1}a(w,1))=f(a),\ \forall a\in 
\overline{A}^{n_0} \},
\]
is well-defined and non-zero by a cocycle calculation
(and similarly for
${\cal H}(\overline{A}^{n_0}//K_0\cap A^{n_0})^W_{gen}$).

Define $D=D_{G/A}$ on $A$ by
\[
\begin{array}{c}
D_{G/A}(\left[ 
\begin{array}{cc}
a & 0 \\ 
0 & a^{-1}
\end{array}
\right] )=\det (1-Ad(\left[ 
\begin{array}{cc}
a & 0 \\ 
0 & a^{-1}
\end{array}
\right] ))_{{\frak g}/{\frak a}} \\ 
=(1-a^2)(1-a^{-2})=-(a-a^{-1})^2,
\end{array}
\]
for $a\in F^{\times}$. Here ${\frak g},{\frak a}$
denote the Lie algebras of $G,A$. 
This may be pulled back to $\overline{A}$
and we have $D(\overline{a})=\Delta(\overline{a})
\delta(\overline{a})^{-1/2}$, for regular
$\overline{a}\in \overline{A}$.
Define the {\bf orbital integral of $f\in
C_c^\infty (\overline{G})_{gen}$ on $\overline{A}_{reg}$}
by 
\begin{equation}
F_f^{\overline{A^{n_0}}}(a)=|D(a)|^{1/2}\int_{\overline{A^{n_0}}\backslash 
\overline{G}}f(x^{-1}ax)\frac{dx}{da},
\label{eqn:2}
\end{equation}
if $a\in \overline{A}$ is regular.
The Haar measure $da$ on $\overline{A^{n_0}}$ will
be choosen later.
Define the {\bf orbital integral of $f\in
C_c^\infty (\overline{G})_{gen}$ on $\overline{A}^{n_0}_{reg}$}
by
\begin{equation}
F_f^{\overline{A}}(a)=n_0|D(a)|^{1/2}
\int_{\overline{A}\backslash \overline{G}}
f(x^{-1}ax)\frac{dx}{da},
\end{equation}
for $a\in \overline{A^{n_0}}\cap \overline{G}_{reg}$.
Since $Cent(a,\overline{G})=\overline{A}$, this is well-defined. 
Here $da$ is the usual product measure on $\overline{A}
=A\times \mu_n$. (The $n_0$ factor is to make Definitions
\ref{defn:transfer} and (\ref{eqn:character}) below
consistent.)

We will show that $f\longmapsto F_f^{\overline{A^{n_0}}}$ defines a 
surjection
\[
C_c^\infty(\overline{G})_{gen}
\rightarrow 
C_c^\infty(\overline{A}^{n_0})^W_{gen}.
\]

\subsection{Dependence on the cocycle}

Let $\beta$ be a cocycle of $G$ having values in $\mu_n$
and let $\overline{G}=\overline{G}_\beta$ denote the
associated covering group.
Note that the metaplectic algebras 
$C_c^\infty(\overline{G})_{gen}$
and ${\cal H}(\overline{G}//K_0)_{gen}$
depend on the cocycle $\beta$
so let us denote this dependence temporarily by
$C_c^\infty(\overline{G})_{gen,\beta}$
and ${\cal H}(\overline{G}//K_0)_{gen,\beta}$,
resp..
If the cocycle is changed to an equivalent one,
say to $\beta'$, where 
\[
\beta(g,h)=\beta'(g,h)s(g)s(h)s(gh)^{-1}
\]
is the cocycle modified by $s$ 
and $s:G\rightarrow \mu_n$ is any function
satisfying $s(1)=1$
(for example, the Kubota splitting
\cite{KP}), then the two algebras are 
isomorphic.

\begin{lemma} 
\label{lemma:isomorphic}
Let $\beta,\beta'$ be as above, let $R$ denote one of 
the algebras $C_c^\infty(\overline{G})_{gen}$ or
${\cal H}(\overline{G}//K_0)_{gen}$,
and let $s:G\rightarrow \mu_n$ satisfy $s(1)=1$.
There is a canonical isomorphism
\[
\phi:R_\beta\cong R_{\beta'},
\]
as algebras, defined by sending 
$f(g,\zeta)=\zeta f(g,1)$ to
$\phi(f)(g,\zeta)=\zeta s(g) f(g,1)$. In other words,
under this mapping, we have
\[
\phi(f_1*_{\beta'} f_2)(g,\zeta)
=(\phi(f_1)*_{\beta}\phi(f_2))(g,\zeta),
\]
where $*_\beta$ denotes the convolution
with respect to the $\beta$ cocycle
\[
(\phi_1*_\beta \phi_2)(g,\zeta)
=\zeta \int_{G}\phi_1(x,1)\phi_2((x^{-1}g,1))
\beta(x,x^{-1}g)^{-1} dx
\]
and, similarly, $*_{\beta'}$ denotes the convolution
with respect to the $\beta'$ cocycle.
\end{lemma}

{\bf proof}: The verification of this
is straightforward.      $\#$

\subsection{The Harish-Chandra transform of a smooth 
function}

The verification of the following lemma, another cocycle 
calculation, is left to the reader.

\begin{lemma}
The centralizer 
$C_{\overline{G}}(\overline{A})$ of $\overline{A}$
in $\overline{G}$ is $\overline{A}^{n_0}$,
where $n_0=n/gcd(2,n)$.
\end{lemma}

\begin{lemma}
(Weyl integration formula \cite{S}, p. 198) For each
$f\in C_c^\infty (G)$
\[
\int_{G}f(g)dg=\sum_{T\in S}\frac
1{|W_T|}\int_T|D(t)|^{1/2}F_f^T(t)dt,
\]
where 
\[
F_f^T(t)=|D(t)|^{1/2}\int_{T\backslash G}f(x^{-1}tx)\frac{dx}{dt},
\]
and where $D(t)=D_{G/T}(t)=\det (1-Ad(t))_{{\frak g}/{\frak t}}$.
Here ${\frak g},\;{\frak t}$ denote the Lie algebras of $G,T$, resp..
\label{lemma:weylint}
\end{lemma}

The following result might be compared with
Flicker's Lemma 1.1 in \cite{F}.

\begin{lemma} (``Vanishing lemma'')
For $f\in C_c^\infty (\overline{G})_{gen}$,
$a\in \overline{A}_{reg}-\overline{A^{n_0}}$,
we have $F_f^{\overline{A^{n_0}}}(a)=0$.
\label{lemma:141}
\end{lemma}

{\bf proof}: Since $dx$ is a Haar measure, we may replace 
$x$ by $a_1x$, where 
$a_1=(\left( 
\begin{array}{cc}
t_1 & 0 \\ 
0 & t_1^{-1}
\end{array}
\right) ,\varsigma _1)\in \overline{A}$ in the integral
(\ref{eqn:2}) defining $F_f^{\overline{A^{n_0}}}(a)$ without
changing its value. Let 
$a=(\left( 
\begin{array}{cc}
t & 0 \\ 
0 & t^{-1}
\end{array}
\right) ,\varsigma )$. Then 
\[
F_f^{\overline{A^{n_0}}}(a)=(t,t_1)_n^2\cdot F_f^{\overline{A^{n_0}}}(a). 
\]
It follows that $F_f^{\overline{A^{n_0}}}(a)=0$ or $(t,t_1)_n^2=1$. 
By the properties of
the Hilbert symbol \cite{N}, if $(t,t_1)_n^2=1$ for all $t_1$ then
$t\in (F^{\times })^{n_0}$. $\# $

\begin{proposition}
 The map $f\longmapsto F_f^{\overline{A^{n_0}}}$\ defines a surjection
$C_c^\infty (\overline{G})_{gen}\rightarrow
C_c^\infty (\overline{A}^{n_0})_{gen}^W$.
\label{prop:main1}
\end{proposition}

\begin{remark}
\begin{itemize}
\item
There's an analogous result for the Schwartz space ${\cal C}$
of ``rapidly decreasing'' functions: 
 The map $f\longmapsto F_f^{\overline{A^{n_0}}}$ 
defines a continuous
surjection ${\cal C}(\overline{G})\rightarrow
{\cal C}(\overline{A^{n_0}})^W$.
Since it shall not be needed, we omit the proof.
\item
The map $f\longmapsto F_f^{\overline{A^{n_0}}}$
is {\it a priori} only defined as a map to functions on
$\overline{A}_{reg}$. The proposition asserts (as we shall
prove below) that this function extends uniquely to
a smooth function on $\overline{A}$.
\end{itemize}
\end{remark}

{\bf proof:} We shall merely sketch the idea.

Define the map $\Phi$ by $\Phi(f)=F_f^{\overline{A^{n_0}}}$.

We first verify surjectivity.

Let $H=L^2(\overline{A}^{n_0})^W$ and let 
$V$ denote the closure in the $L^2$-topology of $H$
of the image of $\Phi:f\longmapsto F_f^{\overline{A^{n_0}}}$ in $H$.

{\bf Claim}: 
If $\Phi$ is not surjective then $V\not= H$.

The proof of claim is straightforward and left to the reader.

If $V\not= H$, we let 
\[
V^\perp =
\{ \phi\in H\ |\ \int_{\overline{A}^{n_0}}
\phi(a)\overline{g(x)}dx=0, \forall g\in V
\}. 
\]

{\bf Claim}: $H=V\oplus V^\perp$.

The proof of claim is straightforward and left to the reader.

Now we complete the proof of surjectivity. 
Let $\phi \in V^{\perp }$. We
must show that $\phi =0$. For any $F\in V$, 
we have 
\[
\int_{\overline{A}^{n_0}}\phi (a)F(a)da=0.
\]
Let 
\[
\tilde \phi (g)=\left\{ 
\begin{array}{cc}
\phi (a),& {\rm \ if\ }g^x=
a\in \overline{A}^{n_0}_{reg}{\rm \ some\ }x\in \overline{G}, \\ 
0,& {\rm otherwise}.
\end{array}
\right. 
\]
Since $\phi$ is $W$-invariant, $\tilde{\phi}$ is well-defined.
Pick $\epsilon >0$ and approximate 
$F$ by $F_f^{\overline{A^{n_0}}}$, for some 
$f\in C_c^\infty (\overline{G})$
such that 
\begin{equation}
\ \ \ \int_{\overline{A}^{n_0}}
|\phi (a)(F(a)-F_f^{\overline{A^{n_0}}}(a))|da<\epsilon .  
\label{eqn:eps}
\end{equation}
By Lemma \ref{lemma:weylint} and 
(\ref{eqn:eps}), we have 
\[
|\int_{\overline{G}}\tilde \phi (g)f(g)dg|<\epsilon /2.
\]
Since $f\in C_c^\infty (\overline{G})$ 
and $\epsilon >0$ were choosen arbitrarily, we
have that $\tilde \phi =0$ a.e. on $\overline{G}$. 
Thus $\phi =0$ a.e. on $A$.
This proves that $V^{\perp }=0$, hence 
$V=H$, hence $\Phi $ is surjective.

For $f\in C_c^\infty (\overline{G})_{gen}$ and
$a\in \overline{A^{n_0}}_{reg}
=\overline{A^{n_0}}\cap \overline{G}_{reg}$, the
function $x\longmapsto f(x^{-1}ax)$ is compactly supported on
$\overline{A}\backslash \overline{G}$ and 
\[
\begin{array}{c}
F_f^{\overline{A^{n_0}}}(a)=|D(a)|^{1/2}
\int_{\overline{A}\backslash \overline{G}}
f(x^{-1}ax)\frac{dx}{da} \\ 
=\delta (a)^{1/2}\int_N\int_{K_0}f(kank^{-1})dkdn \\ 
=\delta (a)^{1/2}\int_Nf^{K_0}(an)dn,
\end{array}
\]
where 
\[
f^{K_0}(x)=\int_{K_0}f(kxk^{-1})dk.
\]

From the claim it 
follows that $F_f^{\overline{A^{n_0}}}(a)$ is
compactly supported on $\overline{A}_{reg}.$

Moreover, since $f$ is locally constant so is $f^{K_0}$. 
Since $f^{K_0}$ is
compactly supported, $\int_Nf^{K_0}(an)dn$ is a locally 
constant function on 
$\overline{A}_{reg}$.

It remains to show that $F_f^{\overline{A^{n_0}}}(a)$ 
extends to a locally constant
function on all of $\overline{A}$. This follows from 
the Shalika germ expansion for covering groups \cite{V}.
 $\# $

The Proposition above implies the following result.

\begin{corollary}
Assume that the cocycle $\beta$ is trivial on $K_0$.
The map $S$ defines an algebra homomorphism
\[
{\cal H}(\overline{G}//K_0)_{gen}
\rightarrow 
{\cal H}(\overline{A}^{n_0}//K_0\cap A^{n_0})^W_{gen}.
\]
\label{cor:main1}
\end{corollary}

\begin{definition} 
Define $N:A\rightarrow \overline{A}$
by $N(a)=(a^{n_0},1)$. We say that
$\overline{f}\in C_c^\infty (\overline{G})_{gen}$
{\bf transfers}\footnote{In other words, the Harish-Chandra
transform of $f$ corresponds with that of $\overline{f}$.}
 to $f\in C_c^\infty(G)$ if
they satisfy
\[
F_{\overline{f}}^{\overline{A}}(N(a))=F_{f}^{A}(a),
\]
for all $a\in A$ such that $N(a)\in \overline{A}_{reg}$,
where the measures have been normalized as above.

A similar definition, for 
$\overline{f}\in {\cal H}(\overline{G}//K_0)_{gen}$
and $f\in {\cal H}(G//K_0)$, is left to the reader.
\label{defn:transfer}
\end{definition}

\section{Representations}

The goal of this section is to describe
the unitary dual of $\overline{G}$.
The general idea in the case of the non-algebraic
group $\overline{G}$ is basically the same as
for the group $G$: using Jacquet functors, one may
show that an irreducible admissible representation 
is either a constituent of some 
representation induced from 
(a supercuspidal representation
of a Levi component of) a proper parabolic subgroup
or else is supercuspidal. As indicted above, the supercuspidal
representations have been described elsewhere.
To capture those irreducible representations 
which are unitary, one using results in
\cite{FK} (for $\overline{GL(2,F)}$), \cite{Mo1},
\cite{Mo2}, and \cite{Ar}. The only cases not falling
into this net are the non-tempered representations
and, if any, the reducible principal series.

For a reductive p-adic group $H$, let
$\Pi (H)$ denote the set of equivalence classes of irreducible
admissible representations of $H$ and let
$\Pi_u(H)$ denote the unitary ones. If $H\subset G$
is a subgroup, let $\Pi_{gen}(\overline{H})$
denote the set of equivalence classes of irreducible
admissible genuine representations $\overline{H}$.

Let $M=A_0$ and $N=N_0$. Each $(\sigma ,W)\in \Pi (\overline{M})$
expends trivially to a representation of
$\overline{B}=\overline{M}N$. Let
$I_{\overline{M}}(\sigma ):\overline{G}\rightarrow Aut(V)$ denote
the {\bf unitarily induced representation}: the representation of
$\overline{G}$ by right translation on 
\begin{eqnarray*}
V &=&\{f:\overline{G}\rightarrow W\;\;{\rm smooth,\ genuine}\;|\; \\
&& 
\begin{array}{c}
(1)\;\;f(bg)=\delta _M(b)^{1/2}\sigma (b)f(g),\;\; 
\;\forall g\in \overline{G},\;b\in \overline{B} \\ 
(2)\;\;{\rm for\;some\; compact\;open\;subgroup}\;K\subset 
\overline{G},\\
\ \ \ f(gk)=f(g),\; \forall k\in K,\;g\in \overline{G}\; \}.
\end{array}
\end{eqnarray*}
Of course, if $\sigma$ is genuine, so is
$I_{\overline{M}}(\sigma )$.


\subsection{Maximal abelian subgroups of $\overline{A}$}

\label{subsec:4.1}
For the purposes of explicitly computing the representations
of $\overline{A}$, we want to determine, as best as we
can, a maximal abelian subgroup of the $n$-fold cover of
$\overline{A}$. This turns out to be hard if $p|n$.
Since the (incomplete) results in the
case $p|n$ will not needed for our final results,
we merely state the result in the case $p$ and
$n$ are relatively prime.

Let
\[
h(x)=\left(
\begin{array}{cc}
x & 0\\
0 & x^{-1}
\end{array}
\right),\ \ \ \ 
x\in F^\times .
\]

\begin{lemma}
A subgroup $\overline{C}\subset \overline{A}$\ is abelian 
if and only
if $(a,a^{\prime })_n^2=1$\ for all $h(a),h(a^{\prime })\in C$.
\end{lemma}

The (easy) proof of this is omitted.

For any $p$-adic field, we have
\[
F^{\times }=\pi ^{{\Bbb Z}}\cdot \mu _{q-1}\cdot U_1, 
\]
a direct product (as abelian groups - see
for example \cite{N}, Proposition 1.1). Here 
\[
U_k=\{x\in {\cal O}_F\ |\ x\equiv 1\  ({\rm mod}\ \pi^k )\}.
\]
If $gcd(p,2n)=1$ then $\mu _n\subset F^{\times }$ implies
$q\equiv 1\;({\rm mod}\;n)$.

\begin{lemma}
\begin{enumerate}
\item[(a)] If $gcd(p,n_0)=1$\ then 
$C=h(\pi ^{{\Bbb Z}}{\cal O}_F^{\times n_0})=
h(\pi ^{{\Bbb Z}}\mu_{q-1}^{n_0}(1+\pi {\cal O}_F))$
and $C=h(\pi ^{{n_0\Bbb Z}}{\cal O}_F^\times )$
are maximal subgroups of 
$F^{\times }$ for which 
$\overline{C}\subset \overline{A}$\ is abelian.

\item[(b)]
If $gcd(p,n)=1$
then $C=h(\pi ^{{\Bbb Z}}(1+\pi {\cal O}_F)\mu _{q-1}^{n_0})$ 
(resp. $C=h(\pi ^{{n_0\Bbb Z}}{\cal O}_F^\times )$)
has index $n_0$ in $F^{\times }$.
\end{enumerate}
\end{lemma}

\begin{remark}
If $n=1$ or $n=2$ then $n_0=1$ and $C=A$.
\end{remark}

The proofs are omitted.

\subsection{Irreducible representations of $\overline{A}$}

\label{subsubsec:4.2.1}

We assume $gcd(p,2n)=1$ throughout this subsection unless otherwise
stated. 

Let $C=h(\pi ^{{n_0\Bbb Z}}{\cal O}_F^\times )$ and 
$A_n=h(\pi ^{{n\Bbb Z}}{\cal O}_F^\times )$. We know that
$\overline{A_n}$ splits, so we may identify
$\overline{A_n}$ with $A_n\times \mu_n$. Moreover,
Clifford-Mackey theory implies that
\[
\begin{array}{ccc}
\Pi(\overline{C})_{gen} &\rightarrow &\Pi(\overline{A})_{gen}\\
\chi & \longmapsto & Ind_{\overline{C}}^{\overline{A}}(\chi)
\end{array}
\]
is a surjection \cite{Mo1}. The conjugates $\chi^{g}$, for 
$g\in \overline{A}/\overline{C}$, are distinct,
and $A/C$ is order $n_0$, so the dimension of each 
$Ind_{\overline{C}}^{\overline{A}}(\chi)$ is $n_0$.

\begin{lemma}
(\cite{Mo2}, \S 1.2)
There is a 2-1 correspondence
\[
\Pi(A_n)\rightarrow \Pi(\overline{A})_{gen}
\]
given by
$\eta_j :\Pi(A_n)\rightarrow \Pi(\overline{C})_{gen}$,
$\iota :\Pi(\overline{C})_{gen}\rightarrow \Pi(\overline{A})_{gen}$,
where 
\[
\iota(\chi)=
Ind_{\overline{C}}^{\overline{A}}(\chi),
\]
and 
\[
\eta_j(\chi)(a,\zeta)=
\left\{
\begin{array}{cc}
\zeta \chi(a), &a\in A_n,\\
\zeta (x,\pi)_n^{n_0}(\pi,\pi)_n^{n_0}\alpha_j\chi(a), 
&a=h(x)\in C-A_n
\end{array},
\right.
\]
where $\alpha_1,\alpha_2$ are the two roots of 
$\alpha^2 =(\pi,\pi)_n^{n_0^2}$.
\end{lemma}

\begin{lemma}
Let $\mu\in \Pi(\overline{C})_{gen}$ and let
$\chi=\chi_\mu =Ind_{\overline{C}}^{\overline{A}}\mu$.
Then 
\begin{equation}
\begin{array}{c}
{\rm tr}\,\chi (a)=\sum_{h\in \overline{A}/\overline{C}}\mu ^h(a)\\
=\left\{
\begin{array}{cc}
n_0\chi(a),& a\in \overline{A^{n_0}},\\
0& a\in \overline{A}-\overline{A^{n_0}}.
\end{array}
\right.
\end{array}
\end{equation}
\label{lemma:chi}
\end{lemma}

{\bf proof}: The first part of this lemma is a special case of the 
more general
Frobenius formula for the trace of an induced representation.
The second part is true, by the orthogonality
of characters, since the character $x\longmapsto (a,x)_n^2$ is non-trivial
if and only if $a\in F^{\times}-F^{\times n_0}$. $\#$


\subsubsection{Corresponding functions on $\overline{A}_{reg}$}

\label{subsubsec:4.2.2}

Define the Fourier transform on 
$\overline{A}$ by
\[
\phi ^{\wedge }(\chi )=\int_{\overline{A}}
{\rm tr}\,\chi (x)\phi (x)dx \\ 
\ \ \ \ \ \ \chi \in \Pi (\overline{A}),
\ \phi\in C_c^\infty(\overline{A}).
\]
If $H$ is an abelian subgroup of $G$ or
of $\overline{G}$ then we define the {\bf Fourier transform}
by
\[
\phi ^{\wedge }(\mu )=\int_{H}
\mu (x)\phi (x)dx \\ 
\ \ \ \ \ \ \mu \in \Pi (H),
\ \phi\in C_c^\infty(H).
\]

We investigate the conditions on $\phi \in C_c^\infty(A)$,
$\overline{\phi}
\in C_c^\infty(\overline{A})_{gen}$
under which
\[
\phi^{\wedge}(\chi)=
\overline{\phi}^{\wedge}(\overline{\chi}),
\]
if $\chi \in \Pi(A)$ and $\overline{\chi}$ are
corresponding representations (in a sense made 
precise below).

\begin{definition} We say that
$\phi \in C_c^\infty(A)$ {\bf corresponds} to
$\overline{\phi}
\in C_c^\infty(\overline{A})_{gen}$
if
\[
\phi^{\wedge}(\chi)=
\overline{\phi}^{\wedge}(\overline{\chi}),
\]
for all $\chi \in \Pi(A_n)$, regarded as an element of
$\Pi(A)$ trivial on $A/A_n$, where $\overline{\chi}=
\iota(\eta_j(\chi))$, for some $j=1,2$.
\end{definition}

\begin{definition} We say that
$f \in C_c^\infty(G)$ {\bf corresponds 
on the hyperbolic set} to
$\overline{f}\in C_c^\infty(\overline{G})_{gen}$
if $F^A\in C_c^\infty(A)$ corresponds to
$F^{\overline{A}}_{\overline{f}}\in C_c^\infty(\overline{A}^{n_0})_{gen}$
in the sense above.
\end{definition}


\subsection{Jacquet functors}

It is remarked in \cite{BD}, \S 2.2
that the arguments of \cite{BZ1}, chapter 2 carry
over to finite central extensions of reductive groups 
over a $p$-adic field (see also \cite{KP}, \S 1.2).
The arguments of \cite{BZ2}, section 2 and
the corresponding sections of \cite{Ca} also
carry over to finite central extensions
of split reductive groups over a $p$-adic field. 
Such results reduce
the determination of the unitary dual of $\overline{G}$
down to classifying the supercuspidal representations (done in \cite{M1}, 
\cite{M2}, \cite{J3} when $gcd(p,n)=1$ and
\cite{J2} for any $p,n$) and the constituents of the
induced representations (see below). 

If $(\pi ,V)$ denotes a $\overline{G}$-module and if 
$P=MN$ denotes a standard parabolic of $G$, let $V_{\overline{N}}$ 
denote the Jacquet module 
\[
V_{\overline{N}}=V/\{\pi (n)v-v\ |\ v\in V,\ n\in \overline{N}\}. 
\]
This sends admissible representations to admissible 
representations (this is
due to Jacquet, see \cite{BZ1}, sections 3.16-3.17
for the $GL(r)$ case).

\begin{lemma}
(Jacquet \cite{Ca}, Theorem 5.2.1)
If $V_{\overline{N}}=0$ for all proper
standard parabolics $P$ of $G$ then $(\pi ,V)$ is supercuspidal.
\end{lemma}

Let ${\cal L}(G)$ denote the set of all standard Levi's 
(with respect to
$A$), up to associates. Therefore, ${\cal L}(G)=\{A,G\}$.

\begin{proposition}
(Jacquet \cite{BZ1}, section 3.19)
If $\pi \in \Pi (\overline{G})$ then
there is a Levi $M\in {\cal L}(G)$ and a supercuspidal
$\sigma \in \Pi (\overline{M})$
such that $\pi $ is a constituent of 
$I_{\overline{M}}(\sigma )$.
\end{proposition}

In particular, every $\pi \in \Pi (\overline{G})$ which is not a
constituent of an induced representation of a supercuspidal 
representation of a
proper Levi occurs discretely in the decomposition 
of the right regular
representation of $\overline{G}$ on $L^2(\overline{G})$.

\subsection{Principal series}

We can investigate whether or not two induced representations 
are equivalent
by using the Jacquet functor. 

Let $\chi ,\chi ^{\prime }\in \Pi (\overline{A})$ and let
$\overline{w}=(w,1)$, for $w\in W$. 
If $\chi ^{\overline{w}}\not
=\chi $ for all $w\in W-\{1\}$ then we call $\chi $ 
{\bf regular}. We say
that $\chi ,\chi ^{\prime }$ are {\bf $W$-conjugate} 
if $\chi ^{\prime}
=\chi ^{\overline{w}}$ for some $w\in W$. The following two 
results were proven for
metaplectic covers of $GL(r,F)$ in \cite{KP}, however they hold in the
present case as well.

\begin{lemma}
(\cite{BZ2}, Corollary 2.13, Theorem 2.9(b))
\label{lemma:bz2}
\begin{enumerate}
\item[(a)]
 Let $\chi \in \Pi (\overline{A})$. The
Jordan-Holder series of 
$I_{\overline{A}}(\chi )_{\overline{N}_0}$ has as its
composition factors 
\[
Ind_{\overline{C}}^{\overline{A}}(\chi ^{\overline{w}}\cdot
\delta _{A}^{1/2}),\ \
w\in W.
\]
If $\chi $ is regular then 
\[
I_{\overline{A}}(\chi )_{\overline{N}_0}=
\oplus _{w\in W}Ind_{\overline{C}}^{\overline{A}}
(\chi^{\overline{w}}\cdot \delta _{A}^{1/2}).
\]

\item[(b)] 
Let $\chi ,\chi ^{\prime }\in \Pi (\overline{A})$.
If $\chi $ is regular then 
\[
dim\ Hom_{\overline{G}}(I_{\overline{A}}(\chi ),
I_{\overline{A}}(\chi ^{\prime }))\leq 1,
\]
with equality if and only if $\chi ^{\prime },\chi $ 
are $W$-conjugate.
\end{enumerate}
\end{lemma}

In other words, distinct $W$-conjugacy classes of
$\chi \in \Pi (\overline{A})$ yield
inequivalent representations.

Suppose that $\pi \in \Pi _u(\overline{G})$. We call 
$\pi $ a (unitary) 
{\bf principal series} representation if $\pi =
I_{\overline{A}}(\chi )$ for some $\chi \in \Pi _u(\overline{A})$. 
These representations are tempered. In case $I_{\overline{A}}(\chi )$
is reducible and $\chi \in \Pi _u(\overline{A})$, we call the
irreducible constituents {\bf reducible principal series}
(or, more precisely, {\bf reducible principal series constituents}).

Let $\chi \in \Pi (\overline{A})$. The induced representation
$I_{\overline{A}}(\chi )$ is in general not irreducible. 
However, we do have the
following

\begin{proposition}
(Moen \cite{Mo2}) 
\begin{enumerate}
\item[(a)]
If $n$ is even and $gcd(p,n)=1$ then 
$I_{\overline{A}}(\chi )$ is
irreducible and unitary for all $\chi \in \Pi _u(\overline{A})$.
\item[(b)]
If $n$ is odd and $gcd(p,n)=1$ then $I_{\overline{A}}(\chi )$ is
irreducible and unitary for all $\chi \in \Pi _u(\overline{A})$ such that 
(a) $\chi =1$ or (b) 
$\chi^{\overline{w_0}}\neq \chi $ where $\overline{w_0}=(w_0 ,1)$. 
If $\chi^{\overline{w_0}}=\chi $ and $\chi\not= 1$ 
then $I_{\overline{A}}(\chi )$ is 
reducible and has two irreducible constituents.
\end{enumerate}
\end{proposition}

In fact, C. Moen \cite{Mo1}
explicitly computes the intertwining operators as
matrices using the Kirillov model when $gcd(p,n)=1$. 

\begin{proposition}
$I_{\overline{A}}(\chi )$ is
irreducible and unitary for all $\chi \in \Pi _u(\overline{A})$
such that 
$\chi^{\overline{w_0}}\neq \chi $ where $\overline{w_0}=(w_0 ,1)$. 
\label{prop:moen}
\end{proposition}

The above result has a direct proof, based on Bruhat theory,
but the following argument deduces the this result from
results in \cite{FK}. 

{\bf proof}: The ``irreducibility if and only if
$\chi^{\overline{w_0}}\neq \chi $'' follows from the
above Lemma \ref{lemma:bz2}(b).

We claim that each representation
$I_{\overline{A}}(\mu )$ as in the proposition
extends to a representation $I'_{\overline{A}}(\mu )$
of $\overline{GL(2,F)}$ (where
$\overline{GL(2,F)}$ is defined by any cocycle of
$GL(2,F)$ extending that of $\overline{SL(2,F)}$
used above). This claim follows from ``Clifford-Mackey theory''
and its proof is omitted (but see Schultz \cite{Sch} and
Gelbart and Piatetski-Shapiro \cite{GPS}
for a special case).

We claim that each such representation
$I'_{\overline{A}}(\mu )$ of $\overline{GL(2,F)}$
is unitary. This claim follows from
\cite{FK}, \S\S 26-27 (which does not assume that
$p$ and $n$ are relatively prime). Let $V_\mu$ denote the
vector space (having a positive-definite Hermitian inner product)
on which the representation $I'_{\overline{A}}(\mu )$ 
of $\overline{GL(2,F)}$ acts.

It follows from these claims that 
$I_{\overline{A}}(\mu )$ acts on a vector space (namely, $V_\mu$)
having a positive-definite Hermitian inner product on it.
$\#$

\subsection{Complementary series}

In this subsection, we shall briefly
review some of the results of Ariturk \cite{Ar}
and use some results of Flicker and Kazhdan \cite{FK}
to generalize them to the $n$-fold cover (\cite{Ar}
assumed $n=3$ and $p>3$).

We call an irreducible unitary representation
$\pi $ a {\bf complementary series} representation if 
$\pi =I_{\overline{A}}(\chi )$ 
for some 
$\chi \in \Pi (\overline{A})-\Pi _u(\overline{A})$. 
These representations are not
tempered.

Let $\mu \in \Pi (\overline{C})$,
$\chi=\chi_\mu =Ind_{\overline{C}}^{\overline{A}}\mu \in
\Pi (\overline{A})$. If $\mu (x)=\mu _0(x)|x|^s$, for some
character $\mu _0$ of finite order and some 
$s\in {\Bbb C}$ then
we write $s=s(\mu )=s(\chi )$.

Let $K(\mu )$ denote the space of locally constant 
functions $f:F\times 
\overline{A}\rightarrow {\Bbb C}$ such that

(i) $f(x,a_1a_2)=\mu (a_1)f(x,a_2)$, 
$a_1\in \overline{C}$, $a_2\in 
\overline{A}$,

(ii) $|x|\chi (\left[ 
\begin{array}{cc}
x & 0 \\ 
0 & x^{-1}
\end{array}
\right] ,1)f(x,a)$ is constant for $|x|$ large.

\noindent
Let $R\subset \overline{A}$ denote a complete set of 
representatives of 
$\overline{A}/\overline{C}$, and let $r$ denote the 
cardinality of $R$.
The elements $f\in K(\mu )$ may be identified with 
the $r$-tuple $%
(f(x,a))_{a\in R}$.

Let $V(\mu )$ denote the space of all locally constant 
functions 
$\varphi :\overline{G}\times \overline{A}\rightarrow 
{\Bbb C}$ such that

(i) $\varphi (g,a_1a_2)=\mu (a_1)\varphi (g,a_2)$, 
$a_1\in \overline{C}$, $a_2\in \overline{A}$,

(ii) $\varphi (a_1ng,a_2)=\delta (a_1)\varphi (g,a_2a_1)$, 
where $a_1\in 
\overline{A}$, $a_2\in \overline{A}$, $n\in N_0$.
\noindent
Here $\delta $ denotes the usual modulus function 
(extended to $\overline{A}$ via the obvious 
pull-back). For $\varphi \in V(\mu )$, define
the map $T$ by
\[
T\varphi (g,a)=\int_F\varphi (w\cdot \left[ 
\begin{array}{cc}
1 & x \\ 
0 & 1
\end{array}
\right] \cdot g,waw^{-1})dx,\;\;Re(s(\mu ))>0\ . 
\]

\begin{lemma} 
(Ariturk)
$T$ intertwines $I_{\overline{A}}(\mu )$ and 
$I_{\overline{A}}(\mu ^w)$.
\end{lemma}

This lemma does not require us to assume $gcd(p,n)=1$.

Let $L(\overline{G},\overline{B})$ denote the space 
of all locally
constant functions $\varphi $ on $\overline{G}$ 
such that 
\[
\varphi ((\left[ 
\begin{array}{cc}
a & * \\ 
0 & a^{-1}
\end{array}
\right] ,\varsigma )\cdot g)=|a|^2\varphi (g)\ .
\]
For $\varphi _1\in V(\mu )$, $\varphi _2\in V(\mu ^w)$, 
the function 
\[
g\longmapsto \int_{\overline{A}/\overline{C}}
\varphi _1(g,a)\varphi
_2(g,a)da 
\]
belongs to $L(\overline{G},\overline{B})$. Therefore, 
\begin{eqnarray*}
&<&\varphi _1,\varphi _2>=\int_{\overline{B}\backslash 
\overline{G}}
\int_{\overline{A}/\overline{C}}
\varphi _1(g,a)\varphi _2(g,a)dadg \\
&=&\int_F\int_{\overline{A}/\overline{C}}
\varphi _1(w^{-1}\cdot \left[ 
\begin{array}{cc}
1 & x \\ 
0 & 1
\end{array}
\right] ,a)\varphi _2(w^{-1}\cdot \left[ 
\begin{array}{cc}
1 & x \\ 
0 & 1
\end{array}
\right] ,a)dad^{\times }x
\end{eqnarray*}
gives a non-degenerate bilinear form on 
$V(\mu )\times V(\mu ^{\overline{w}})$.

\begin{lemma}
(Ariturk)
$I_{\overline{A}}(\mu ^{\overline{w}})$ is the contragredient of 
$I_{\overline{A}}(\mu )$.
\end{lemma}

This lemma does not require us to assume $gcd(p,n)=1$.

For $f\in K(\mu)$, define the Fourier transform of
$f$ by
\[
f^{\wedge }(x,a)
=\int_{F^\times}
f(y,a)\psi(xy)dy,
\]
where $\psi$ is a fixed additive character of $F$.

\begin{lemma}
(Ariturk)
Assume $gcd(p,n)=1$.
For $\varphi _1,$ $\varphi _2\in V(\mu )$, 
$\mu (x)=|x|^s$, we have 
\[
<\varphi _1,T\varphi _2>=
\int_F\int_{\overline{A}/\overline{C}}
f_1^{\wedge }(x,a)\overline{(Jf_2^{\wedge })(x,a)}dadx,
\]
where $J=J_\mu $ is a linear transformation on 
$K(\mu )^{\wedge }$ and 
\[
f_i(x,a)=\varphi _i(w^{-1}\cdot \left[ 
\begin{array}{cc}
1 & x \\ 
0 & 1
\end{array}
\right] ,a),\;\;i=1,2.
\]
\label{lemma:Jdef}
\end{lemma}

We may identify the map $J=J_\mu$ defined in the above lemma
with an $r\times r$ matrix which we still denote by $J$.

\begin{lemma}
(Langlands, Ariturk) Assume $gcd(p,2n)=1$.
If $0\leq Re(s(\mu ))\leq 1/n$ and $|Im(s(\mu ))|\leq
\pi /n\ln (q)$ then the image of $J_\mu $ is an 
irreducible representation of $\overline{G}$.
\end{lemma}

\begin{proposition}
\label{prop:fk}
(Flicker-Kazhdan)
If $0< s(\mu )< 1/n$ then $I_{\overline{A}}(\mu )$ is a 
unitarizable representation of $\overline{G}$.
\end{proposition}

{\bf proof}: 
Each such representation
$I_{\overline{A}}(\mu )$ extends to a
representation $I'_{\overline{A}}(\mu )$
of $\overline{GL(2,F)}$ (where
$\overline{GL(2,F)}$ is defined by any cocycle of
$GL(2,F)$ extending that of $\overline{SL(2,F)}$
used above).
We claim that each such representation
$I'_{\overline{A}}(\mu )$ of $\overline{GL(2,F)}$
is unitary. This claim follows from
\cite{FK}, \S\S 26-27 (which does not assume that
$p$ and $n$ are relatively prime). Let $V_\mu$ denote the
vector space (having a positive-definite Hermitian inner product)
on which the representation $I'_{\overline{A}}(\mu )$ 
of $\overline{GL(2,F)}$ acts.

It follows from these claims that 
$I_{\overline{A}}(\mu )$ acts on a vector space (namely, $V_\mu$)
having a positive-definite Hermitian inner product on it.
$\#$

\begin{corollary} 
If $0<s<1/n$ and $\mu (x)=|x|^s$ then 
$<\varphi _1,T\varphi _2>$ is a positive 
definite form. In
particular, $I_{\overline{A}}(\mu )$ is unitary in this range.
If, in addition, $gcd(p,2n)=1$ then $J_\mu $ is 
a positive definite matrix.
\end{corollary}

\begin{remark}
In the case $n=2$, this follows from \cite{GS}, 
Proposition 1 following the argument of \cite{Ar}. 
See also \cite{Mo1} if $p>2$ and \cite{G}, \S 5.4.
In the case $n=3$, $p >3 $, this was proven in \cite{Ar}. 

Though this corollary is not needed for the
character relations in \S 5, it does indicate that some of
the results of
Ariturk \cite{Ar} generalize to the $n$-fold cover case
without any condition on the prime $p$.
\end{remark}

{\bf proof}: First, recall the following form 
due to Jacquet, but see \cite{C}, page 118) of 
Schur's lemma holds: If $(\pi,V)$ is an irreducible smooth 
representation of $\overline{G}$ then every
intertwining map $T:V\rightarrow V$ is a scalar.

Now we claim: If $(\pi,V)$ is an irreducible smooth 
representation of $\overline{G}$ then every
there is at most one $\overline{G}$-invariant
inner product on $V$. This is a well-known (and not hard to prove)
consequence of Schur's lemma.

The previous proposition says that 
the representation space of $I_{\overline{A}}(\mu )$ 
has a $\overline{G}$-invariant, positive-definite,
Hermitian inner product. It follows
$<\varphi _1,T\varphi _2>$ is positive-definite.
This implies $J_\mu$ is also positive-definite, as claimed
$0<s<1/n$ and $\mu (x)=|x|^s$.
$\#$

\subsection{The special and the ``trash'' representations}
\label{subsec:special}

For the origin of the term ``trash'' representation,
see \cite{G}. 

In the notation above, we have the following

\begin{proposition} 
Let  $s=1/n$ and $\mu (x)=|x|^s$.
\begin{itemize}
\item[(a)]
The irreducible subrepresentation
of $I_{\overline{A}}(\mu)$
(if $gcd(p,2n)=1$, the kernel of $J_\mu $)
is ``the special representation'' $\pi _{sp}$. 
It is tempered and square-integrable
(hence unitary). If $gcd(p,n)=1$ then it
also contains an Iwahori fixed vector. 
\item [(b)]
If $n>1$ then the irreducible quotient
of $I_{\overline{A}}(\mu)$
(if $gcd(p,2n)=1$, image of $J_\mu $) is 
an infinite-dimensional, non-tempered representation $\pi _{nt}$.
If $gcd(p,n)=1$ then it is also spherical.
\end{itemize}
\end{proposition}

\begin{remark}
In the case $n=3$ and $p>3$, this 
proposition follows from \cite{Ar}. 
In case $n=2$, most of the statements
are proven in \cite{GS}. 

\end{remark}

{\bf proof}: These statements follow easily from results of Flicker
\cite{F} and Flicker-Kazhdan \cite{FK}, using an
inversion formula relating the character of a representation
of $\overline{G}$ to values of the character of
$\overline{GL(2,F)}$ which it extends to (see page 39 of \cite{Sch}).
$\#$

\begin{proposition} (Kazhdan-Patterson)
$\pi _{nt}$ is unitary.
\end{proposition}

\begin{remark}
This is true if $n=2$ (\cite{GS}, Theorem 2) or
if $n=3,\ p\not= 3$ (\cite{Ar}, Theorem 5.4).
\end{remark}

{\bf proof}\footnote{I thank the referee for this proof.}:
This follows from Kazhdan-Patterson
\cite{KP}, Theorem II.2.1. $\#$

\section{Character relations}

In this section, we use the above results to determine
a ``natural'' correspondence between the (unramified)
principal series of $\overline{G}$ and (unramified)
principal series of $G$. We assume throughout this section 
that $gcd(p,2n)=1$, since we shall require
some of the results in \S \ref{subsec:4.1} above.

We show in this section that if

\begin{itemize}
\item
the image of $\overline{f}$ under
the Harish-Chandra transform
\[
C_c^\infty(\overline{G})_{gen}
\rightarrow 
C_c^\infty(\overline{A}^{n_0})^W_{gen}
\]
``matches up with'' (in a way we will make explicit below)
the image of $f$ under the Harish-Chandra transform
\[
C_c^\infty(G) \rightarrow C_c^\infty(A)^W,
\]
and 

\item 
the principal series representation 
$\overline{\pi}=I_{\overline{A}}(\chi)$
``corresponds with'' (in a way we will make explicit below)
 the principal series representation $\pi=I_{A}(\mu)$,
\end{itemize}
then we have a character relation of the form
\[
{\rm tr}\,\pi(f)={\rm tr}\,\overline{\pi}(\overline{f}).
\]

We also give an analog of this result for spherical functions.

\subsection{Character calculations for principal series}

Let $A$ denote the diagonal subgroup of $G$.
Recall $W=N_{G}(A)/A$ denotes the Weyl group of $A$. 
If $x\in \overline{G}$ is regular then let
$\overline{G}_t=Cent(t,\overline{G})$.

\begin{lemma}
Let $h$ denote an ``anti-genuine'' class function on $\overline{G}$
and suppose that $h$ satisfies the property that, for all
$f\in C_c^\infty(\overline{G})_{gen}$, we have
\[
\int_{\overline{G}}
h(g)f(g)dg=
c\int_{\overline{A}}\Delta(a)^{1/2}h(a)
\int_{\overline{G}_a\backslash \overline{G}}
f(g^{-1}ag){dg\over da}da,
\]
for some constant $c$. Then the following holds:

(1) $c={1\over |W_A|}$,

(2) if $g\in \overline{G}$ is regular then
$h(g)=0$ unless $g$ is conjugate to an element of 
$\overline{A}$.
\label{lemma:Weyl}
\end{lemma}

{\bf proof}: First, note that the integral
$\int_{\overline{G}_a\backslash \overline{G}}
f(g^{-1}ag){dg\over da}da$ is zero unless
$a\in \overline{A^{n_0}}$, in which case
it is equal to $\int_{\overline{A}\backslash \overline{G}}
f(g^{-1}ag){dg\over da}da$. Lemma \ref{lemma:Weyl}
above is now a corollary of the non-metaplectic version, 
which is equivalent to Lemma
\ref{lemma:weylint}. $\#$


\begin{proposition}
Assume $gcd(p,2n)=1$. Assume $\chi =\chi_\mu=
Ind_{\overline{C}}^{\overline{A}}\mu \in 
\Pi (\overline{A})$ and suppose that $\chi \neq \chi ^w$. 
The character of the
irreducible representation $\pi =I_{\overline{A}}(\chi )$
is given on the regular set by 
\[
\Theta _\pi (x)=\left\{ 
\begin{array}{c}
\frac{{\rm tr}\,\chi (a)+{\rm tr}\,(\chi (a)^{-1})}{|a-a^{-1}|},\;\;
x=a^g,\;%{\rm some}
\;g\in \overline{G},\;a\in \overline{A}_{reg}^{n_0}, \\ 
0,\;\;\;\;\;\;\;\;\;\;\;\;\;\;\;\;\;\;\;\;\;
{\rm \;if\;}\;
x\in \overline{G}_{reg}-(\overline{A}^{n_0})^{\overline{G}},
\end{array}
\right. 
\]
where 
\[
(\overline{A}^{n_0})^{\overline{G}}=\{a^g\;|\;
a\in \overline{A}^{n_0},\;g\in 
\overline{G}\}.\ 
\]
Moreover, ${\rm tr}\,\chi(a)=0$ unless 
$a\in \overline{A}^{n_0}$.
\label{prop:trace}
\end{proposition}

{\bf proof}: The proof, which follows Theorem 5.14
in \cite{G}, is omitted. $\#$

Proposition \ref{prop:trace}, Lemma \ref{lemma:chi}, and 
Lemma \ref{lemma:141}
imply that, for $f\in C_c^\infty (\overline{G})_{gen}$, 
\begin{equation}
\begin{array}{c}
{\rm tr}\,\pi (f)=\int_{\overline{G}}\Theta _\pi (x)f(x)dx \\ 
=\frac 12\int_{A^{n_0}}({\rm tr}\,\chi (a,1)+{\rm tr}\,(\chi
(a,1)^{-1}))F_f^{\overline{A}}(a,1)da \\ 
={1\over 2}\sum_{h\in \overline{A}/\overline{C}}
(F_f^{\overline{A}})^{\wedge }(\mu ^h)
+(F_f^{\overline{A}})^{\wedge }((\mu^{\overline{w_0}}) ^h)\\
={1\over 2}|A/C|[(F_f^{\overline{A}})^{\wedge }(\mu )
+(F_f^{\overline{A}})^{\wedge }(\mu^{\overline{w_0}})].
\end{array}
\label{eqn:character}
\end{equation}

\subsection{Matching principal series characters}

Note that since $F^{\overline{A}}_f$ (for
$f\in C_c^\infty(\overline{G})_{gen}$)
vanishes off $\overline{A}^{n_0}$,
the last Fourier transform in (\ref{eqn:character}),
 {\it a priori} an integral over $\overline{C}$,
is in fact an integral over $\overline{A^{n_0}}$.
This, with Definition \ref{defn:transfer},
proves the following consequence of the above proposition.

\begin{theorem}
Assume $gcd(p,2n)=1$.
If $\chi=\chi_\mu=Ind_{\overline{C}}^{\overline{A}}\mu \in 
\Pi (\overline{A})$, $\chi \neq \chi ^{\overline{w_0}}$, and if 
$\overline{\pi} =I_{\overline{A}}(\chi)$
then the character of $\overline{\pi}$ is related to that of
$\pi =I_{A}(\mu^{n_0} )$ by
\[
{\rm tr}\,\overline{\pi}(\overline{f})
={\rm tr}\,\pi(f),
\]
provided $\overline{f}\in C_c^\infty(\overline{G})_{gen}$ 
transfers to $f\in C_c^\infty(G)$.
\label{thrm:main1}
\end{theorem}
 
We now give an analogous result for the 
spherical functions and the unramified principal series.

\begin{theorem}
Assume $gcd(p,2n)=1$.
Assume that the cocycle $\beta$ defining $\overline{G}$
is trivial on $K_0$. If $\chi=\chi_\mu=
Ind_{\overline{A}^{n_0}}^{\overline{A}}\mu \in 
\Pi (\overline{A})$ is unramified, $\chi \neq \chi ^{\overline{w_0}}$, and if 
$\overline{\pi} =I_{\overline{A}}(\chi)$
then the character of $\overline{\pi}$ is related to that of
$\pi =I_{A}(\mu^{n_0} )$ by
\[
{\rm tr}\,\overline{\pi}(\overline{f})
={\rm tr}\,\pi(f),
\]
provided $\overline{f}\in 
{\cal H}(\overline{G}//K_0)_{gen}$
transfers to $f\in {\cal H}(G//K_0)$.
\label{thrm:main2}
\end{theorem}

{\bf proof}: Since $I_{\overline{A}}(\chi_\mu)$ is unramified
if and only if $\mu$ is, this is a corollary of the previous theorem. 
$\#$



\begin{thebibliography}{99}

\bibitem[Ar]{Ar} Ariturk, ``On the composition series of principal series
representations of a three-fold covering group of $SL(2,K)$'', Nagoya Math
J. \underline{77}(1980)177-196

\bibitem[BD]{BD} J. Bernstein, ``Le centre de Bernstein'', 
(notes by P. Deligne), 
in {\bf Repr\'esentations des groupes r\'eductifs 
sur un corp local}, Hermann, Paris, 1984

\bibitem[BZ1]{BZ1} J. Bernstein and A. Zelevinski, ``Representations of 
the group $GL(n,F)$,
where $F$ is a non-archimedian local field'', Russian Math Surveys
\underline{31}(1976)1-68

\bibitem[BZ2]{BZ2} ----, ``Induced representations of reductive p-adic 
groups, I'', Ann.
scient. \'Ec. norm. Sup. \underline{10}(1977)441-472

%\bibitem[Bl]{Bl} C. Blondel, ``Les repr\'esentations supercuspidales 
%des groupes
%m\'etaplectiques sur GL(2) et leurs caract\'eres'',
%Bull de la Soc. Math. France, Memoir 18, 1985

%\bibitem[Bo]{Bo} A. Borel, et al, \underline{Seminar in 
%algebraic groups and
%related finite groups}, Springer LNM \underline{131}, 1970

%\bibitem[BH]{BH} D. Bump and J. Hoffstein, ``On Shimura's 
%correspondence'', 
%Duke Math. J. \underline{55}(1987)661-691

\bibitem[C]{C} P. Cartier, ``Representations of p-adic groups: a survey'',
in {\bf Automorphic forms, representations, and L-functions},
vol I, ed. A. Borel, W. Casselman, AMS, Providence, RI, 1979

\bibitem[Ca]{Ca} W. Casselman, ``Introduction to the theory of admissible
representations on reductive p-adic groups'', preprint

\bibitem[F]{F} Y. Flicker, ``Automorphic forms on covering groups of 
$GL(2)$'', Inv. Math. \underline{57}(1980)119-182

\bibitem[FK]{FK} ------ and D. Kazhdan, ``Metaplectic 
correspondence'', Publ. Math. IHES \underline{64}(1986)53-110

\bibitem[G]{G} S. Gelbart, {\bf Weil's representation and the 
spectrum of the metaplectic group},
SLN Math \underline{530}, 1976

%\bibitem[GK]{GK} ------ and T. Knapp, ``L-indistinguishability and R 
%groups for the special linear group'', Adv. Math. 43(1982)101-121

\bibitem[GPS]{GPS} ------ and I. Piatetski-Shapiro, ``Some remarks 
on metaplectic cusp
forms and the correspondences of Shimura and Waldspurger'', Israel 
J. Math.
\underline{44}(1983)97-126

\bibitem[GS]{GS} ------ and P. Sally, ``Intertwining operators and 
automorphic
froms on the metaplectic group'', Proc. Nat. Acad. Sci.
\underline{72}(1975)1406-1410

\bibitem[J1]{J1} D. Joyner, ``The metaplectic analog of 
Kazhdan's endoscopic lift'',
Israel J. Math. \underline{66}(1988)113-154

\bibitem[J2]{J2} ------, ``A correspondence for the generalized 
Hecke algebra of
the metaplectic cover $\overline{SL(2,F)}$, $F$ p-adic,'' 
New York J. of Math. \underline{4}(1998) 223-235.

\bibitem[J3]{J3} ------, ``A correspondence for the supercuspidal
representations of $\overline{SL_2(F)}$, $F$ p-adic'', 
to appear in Archiv der Math.

\bibitem[KP]{KP} D. Kazhdan and  S. Patterson, 
``Metaplectic forms'', 
Publ. Math. IHES \underline{59}(1984)35-142

\bibitem[K]{K} T. Kubota, ``Topological covering of $SL(2)$ over a 
local field'', J.
Math. Soc. Japan \underline{19} (1967)114-121

%\bibitem[K]{K} P. Kutzko,
%``Character formulas for supercuspidal representations of
%$GL(\ell)$, $\ell$ a prime'', Amer J Math 
%\underline{109}(1987)201-222

%\bibitem[L]{L} R. Langlands, \underline{Base} \underline{change} 
%\underline{for $GL(2)$},
%Annals of Math. Studies, Princeton

\bibitem[M1]{M1} D. Manderscheid, ``On the supercuspidal 
representations of
$SL(2)$ and its 2-fold cover, I'', Math. Ann. 
\underbar{266}(1987)287-295

\bibitem[M2]{M2} ------, ``On the supercuspidal representations of
$SL(2)$ and its 2-fold cover, II'', Math. Ann.
\underbar{266}(1987)297-305

%\bibitem[M3]{M3} ------, ``Supercuspidal representations and the
%theta correspondence'',  J. Algebra \underbar{151}(1992), 
%no. 2, 375--407

\bibitem[Mo1]{Mo1} C. Moen, ``Intertwining operators for covering groups 
of $SL(2)$ over
a local field'', preprint

\bibitem[Mo2]{Mo2} ------, ``Irreducibiliy of unitary principal series 
representations
for covering groups of $SL(2)$'', Pac. J. Math. 135(1988)89-110

\bibitem[N]{N} J. Neukirch, {\bf Class field theory},
Springer-Verlag, 1986

\bibitem[Sch]{Sch} J. Schultz, ``Lifting of characters of 
$\overline{SL_2(F)}$ and
$SO_{2,1}(F)$, for $F$ nonarchimedean'',
PhD thesis (supervised by J. Adams),
Univ. Maryland at College Park, 1998

\bibitem[S]{S} A. Silberger, {\bf Introduction to
harmonic analysis on reductive p-adic groups}, 
Princeton Univ. Press, Princeton, 1979

\bibitem[V]{V} M.-F. Vign\'eras, ``Caract\'erisation des integrales 
orbitales sur un
groupe r\'eductif p-adique'', J. Fac. Sci. Univ. Tokyo 
\underline{28}(1981)945-961

\end{thebibliography}

\vskip 1in
\noindent
Mathematics Dept.\\
Chauvenet Hall\\
572C Holloway Road\\
U.S. Naval Academy\\
Annapolis, MD 21402\\
\verb+wdj@gwmail.usna.edu+
\end{document}




